\documentclass[11pt,a4paper]{moderncv}
\moderncvstyle{classic}
\moderncvcolor{black}
\definecolor{firstnamecolor}{rgb}{0,0,0}

\usepackage[utf8]{inputenc}
\usepackage[T1]{fontenc}
\usepackage[scale=0.85]{geometry}
\usepackage[dvipsnames]{xcolor}
\definecolor{PaleBlue}{RGB}{76, 151, 215}
\usepackage[colorlinks=true, urlcolor=PaleBlue, linkcolor=PaleBlue]{hyperref}
\usepackage{microtype}

\renewcommand{\rmdefault}{ppl}
\renewcommand{\sfdefault}{phv}
\renewcommand*{\namefont}{\fontsize{26}{18}\mdseries}
\microtypesetup{expansion=false}

\input{../../secrets/personal_info_dk.tex}

\title{Curriculum Vitae}

\begin{document}
\makecvtitle

\section{Profile}
Quantitative researcher and data scientist trained in engineering physics, statistical learning, economics, and applied mathematics. I work across probabilistic modelling, optimisation, and Python-based machine-learning systems.

\section{Education}

\cventry{2015--2020}{Royal Institute of Technology (KTH), Stockholm}{MSc Statistical Learning and Data Analytics; BSc Engineering Physics}{GPA: 5.0 / 5.0}{}{Thesis: \href{https://kth.diva-portal.org/smash/get/diva2:1431327/FULLTEXT02.pdf}{See thesis}. \newline Relevant coursework: Quantum Physics; Mathematical Methods in Physics; Probability Theory.}

\cventry{2017--2022}{Stockholm University, Stockholm}{Supplementary Bachelor's Degree in Economics}{}{}{Complemented the engineering background with economic theory and quantitative methods.}

\cventry{2019}{Swiss Federal Institute of Technology (ETH), Zurich}{Exchange, Applied Mathematics}{}{}{}

\section{Experience}

\cventry{2024--Present}{Full-Stack Data Scientist, Project Manager}{Signum Life Science}{Copenhagen, Denmark}{}{
\begin{itemize}
\item Maintain and extend a pharmaceutical price forecasting platform and the Kwarts next-best-action engine.
\item Own and extend Kwarts with explainable AI features, and use embedding-based representations for downstream similarity search and matching workflows.
\item Organise internal knowledge-sharing, including a week-long hackathon with up to 40 participants.
\end{itemize}
}

\cventry{2022--2024}{Principal Data Scientist, Consultant}{Valcon}{Copenhagen, Denmark}{}{
\begin{itemize}
\item Led data science projects across pharmaceuticals, energy, IoT, and manufacturing.
\item Implemented an end-to-end production pipeline for a real-time deep neural network.
\item Presented complex technical topics to non-technical stakeholders, mentored junior colleagues, and led xAI and Git training sessions.
\item Received Valcon's People's Choice Award for exemplary work, leadership, team spirit, and scientific curiosity.
\end{itemize}
}

\cventry{2019--2021}{Data Scientist, Researcher}{RaySearch Laboratories}{Stockholm, Sweden}{}{
\begin{itemize}
\item Developed pipelines for extracting latent anatomical features from segmented 3D CT scans using autoencoders.
\item Predicted dose-volume histograms and spatial dose distributions using sparse Gaussian processes.
\item Redesigned lexicographic optimisation algorithms and collaborated with medical physicists and clinicians on practical workflow fit.
\end{itemize}
}

\section{Technical Skills}
\cvitem{Programming}{Python, SQL, C++, Matlab}
\cvitem{Methods}{Deep learning, autoencoders, CNNs, sparse Gaussian processes, optimisation, explainable AI, time series forecasting}
\cvitem{Tools}{PyTorch, TensorFlow, Scikit-learn, Docker, FastAPI, Databricks, AWS, Qdrant}
\cvitem{Languages}{English (native), Swedish (native), Danish (intermediate)}

\end{document}
